\documentclass[12pt]{article}
\usepackage[utf8]{inputenc}
\usepackage[T1]{fontenc}
\usepackage[margin=1in]{geometry}
\usepackage{setspace}
\usepackage{xcolor}
\usepackage[colorlinks=true, urlcolor=blue!60!black, linkcolor=blue!60!black]{hyperref}
\setlength{\parindent}{2em}
\setlength{\parskip}{1.25em}
\renewcommand{\baselinestretch}{1.0}
\begin{document}
\begin{center}
{\LARGE \textit{Katzenbach v. McClung}, 379 U.S. 294 (1964)\par}
\vspace{0.8em}
{\large Case Brief --- Four-Page Edition\par}
\end{center}
\vspace{1em}

\section{Memory Jogger}

Unanimously, the Court held that Congress may use the Commerce Clause to ban racial discrimination at a purely local restaurant, Ollie's Barbecue, so long as a substantial portion of its food has moved in interstate commerce and Congress had a rational basis to conclude that such discrimination, in the aggregate, burdens interstate trade (\href{https://www.courtlistener.com/opinion/106951/katzenbach-v-mcclung/}{CourtListener}).

\section{Detailed Case Facts}

Ollie McClung, Sr., and Ollie McClung, Jr., ran Ollie's Barbecue in Birmingham, Alabama, a 220-seat family restaurant on a state highway 11 blocks from an interstate. It catered to local families and white-collar workers and, since 1927, had refused to serve African Americans in its dining room, offering only take-out. In the year before the Civil Rights Act of 1964, it bought about \$150,000 of food locally, of which \$69,683 (46\%) was meat a local supplier had procured from outside Alabama.

Title II guaranteed full and equal enjoyment of public accommodations without discrimination based on race, color, religion, or national origin. It covered a restaurant if ``it serves or offers to serve interstate travelers or a substantial portion of the food which it serves . . . has moved in commerce.'' Unlike the companion case \textit{Heart of Atlanta Motel v. United States}, the government made no claim that interstate travelers ate at Ollie's. The sole issue was whether Congress could reach a local restaurant because of its interstate food.

\section{Procedural History}

The McClungs sued Acting Attorney General Nicholas Katzenbach to enjoin enforcement. A three-judge district court held Title II unconstitutional as applied, reasoning that Congress had legislated a conclusive presumption without a demonstrable connection between out-of-state food and restaurant discrimination. Justice Black stayed the injunction; on direct appeal, argued with \textit{Heart of Atlanta}, the Supreme Court reversed on December 14, 1964.

\section{Judicial Votes}

Reversed, \textbf{9--0}.
\begin{itemize}
\item \textbf{Majority:} Justice Clark, joined by Chief Justice Warren and Justices Harlan, Brennan, Stewart, White, Douglas, and Goldberg.
\item \textbf{Concurrences:} Justices Black, Douglas, and Goldberg, each filed in \textit{Heart of Atlanta Motel} and applicable to both cases.
\item \textbf{Dissents:} None.
\end{itemize}

\section{Holding}

Title II is constitutional under the Commerce and Necessary and Proper Clauses as applied to a local restaurant serving food a substantial portion of which has moved in interstate commerce. Congress may regulate local commercial establishments when it has a \textit{rational basis} for finding that the regulated class of activities, in the aggregate, substantially affects interstate commerce and that its scheme is a reasonable means of protecting commerce. Neither formal findings nor case-by-case proof is required. The Court treated the complaint as a declaratory judgment action to settle the question promptly.

\section{Analysis of Opinions}

\textbf{Majority (Clark).} \textit{Heart of Atlanta} disposed of the other constitutional claims and the \textit{Civil Rights Cases}. Congressional hearings showed four links between restaurant segregation and commerce: depressed spending by African Americans meant restaurants bought less out-of-state food; discrimination restricted the market and depressed business generally; it obstructed interstate travel, ``for one can hardly travel without eating''; and it deterred new businesses and skilled workers from relocating. Under \textit{Wickard v. Filburn}, Ollie's trivial contribution, taken together with many others similarly situated, was far from trivial. The ``come to rest'' cases govern state taxation, not federal regulation. Under \textit{Darby} and \textit{Carolene Products}, once a rational basis appears, ``our investigation is at an end.''

\textbf{Black (concurring).} Forty-six percent was plainly ``substantial.'' But whether local operations affect commerce is ``ultimately a judicial, rather than a legislative, question,'' and an isolated lunchroom buying almost everything locally might lie beyond Congress's reach.

\textbf{Douglas (concurring).} The right to be free of discrimination ``occupies a more protected position in our constitutional system than does the movement of cattle, fruit, steel and coal across state lines''; he would rest Title II on \S~5 of the Fourteenth Amendment, which would have a ``more settling effect.''

\textbf{Goldberg (concurring).} The Act's primary purpose is ``the vindication of human dignity, and not mere economics''; Congress had power under both the Commerce Clause and \S~5.

\section{Examples: Future Applications}

\begin{enumerate}
\item \textbf{Same-Side (Valid): Grocery Stores and Coffee Shops.} A ban on discrimination by stores buying at least 25\% of their inventory out of state reaches commercial activity with an express jurisdictional element.
\item \textbf{Same-Side (Valid): Compounding Pharmacy.} A comprehensive drug-market statute reaches a local pharmacy using out-of-state ingredients, consistent with \textit{Gonzales v. Raich}.
\item \textbf{Opposite-Side (Invalid): Knives Near Parks.} A ban on possessing a knife near a public park, with no jurisdictional element, regulates noneconomic conduct and fails under \textit{Lopez}.
\item \textbf{Opposite-Side (Invalid): Compelled Restaurant Meals.} A mandate to buy two restaurant meals a month fails under \textit{NFIB v. Sebelius}, which bars compelling inactive individuals into commerce.
\item \textbf{Fence-Sitter: The Farm-to-Table Diner.} A rural diner sources 99\% of its ingredients locally, using only out-of-state salt and a stove. It is commercial, but Black warned about the ``isolated and remote lunchroom,'' and Title II requires a ``substantial portion'' of interstate food.
\end{enumerate}

\section{Critique}

\textbf{Originalist critique.} Richard Epstein, Raoul Berger, and Justice Thomas view \textit{McClung} as the high-water mark of an unmoored Commerce Clause: ``Commerce . . . among the several States'' originally meant interstate trade, not local retail service after goods come to rest, and rational-basis deference leaves no enforceable boundary. \textit{Lopez} and \textit{Morrison} preserved \textit{McClung} because restaurants are economic, while rejecting its rhetoric for noneconomic conduct.

\textbf{Progressive critique.} Civil-rights scholars applaud the result but, like Douglas and Goldberg, regret resting equal citizenship on interstate meat shipments rather than \S~5, which left later statutes such as the Violence Against Women Act exposed in \textit{Morrison}. Defenders answer that the commerce ground secured unanimity and rapid compliance.

\section{Key Quotations}

\begin{quote}
``Of course, the mere fact that Congress has said when particular activity shall be deemed to affect commerce does not preclude further examination by this Court. But where we find that the legislators, in light of the facts and testimony before them, have a rational basis for finding a chosen regulatory scheme necessary to the protection of commerce, our investigation is at an end.''
\end{quote}
\noindent --- Justice Clark, for the Court, 379 U.S. at 303--304.

\begin{quote}
``The power of Congress in this field is broad and sweeping; where it keeps within its sphere and violates no express constitutional limitation it has been the rule of this Court, going back almost to the founding days of the Republic, not to interfere.''
\end{quote}
\noindent --- Justice Clark, for the Court, 379 U.S. at 305.

\begin{quote}
``We believe that this testimony afforded ample basis for the conclusion that established restaurants in such areas sold less interstate goods because of the discrimination, that interstate travel was obstructed directly by it, that business in general suffered, and that many new businesses refrained from establishing there as a result of it.''
\end{quote}
\noindent --- Justice Clark, for the Court, 379 U.S. at 300.

\end{document}
